\section{Several minimizers}\label{sec:optsets}
The growth analysis of Section~\ref{sec:growth} localizes every retained
coordinate near one minimizer. When the optimal set $\mathcal S$ is not a
single point, this localization fails, although Theorem~\ref{thm:transfer}
shows that exactness itself costs only $\poly(I)+\max\{0,\log_2(1/g_S)\}$
bits of accuracy under set growth with constant $g_S$. This section treats three
situations. With finitely many minimizers, unions of uniform cells give grids
with $O(r\sqrt{n\kappa_S})$ nodes per coordinate, where $r$ is the largest
number of optimal values of one coordinate
(Section~\ref{sec:exact-nonunique}). For coordinatewise concave quadratics,
the decomposition describes the entire optimal set
(Section~\ref{sec:endpointset}). For continuous box quadratics certified by a
diagonal Lagrangian multiplier, an exact optimizer and an exact description
of $\mathcal S$ are found in time $f'(p,\kappa_S)\poly(I)$ for a computable
function $f'$ (Section~\ref{sec:diagcert}). Throughout, $F=\frac12x^\T Hx+b^\T x+c$ is a
rational quadratic on a mixed box, and set growth, $g_S$ and
$\kappa_S=\kappa(L,g_S)$ are as in Definition~\ref{def:growth}. Some $g_S>0$
always exists (Remark~\ref{rem:setgrowth}), but its size is not controlled,
and no algorithm below computes it.

\subsection{One center cannot serve several minimizers}\label{sec:exact-nonunique}

Theorem~\ref{thm:transfer} reduces exact output to reaching accuracy
$\varepsilon_S$. CT (Algorithm~\ref{alg:ct}) does not reach it at a rate
controlled by $p$ and $\kappa_S$ when the minimizer is not unique, even with
two isolated minimizers: its grids are graded around one center, and coarse
cells near a distant minimizer keep the corrected minimum low.

\begin{proposition}[One center cannot certify two distant minimizers]
\label{prop:twocenters}
For $M>0$ let
\[
F_M(x,z)=x^2-2xz+Mz\qquad\text{on }X=[0,M]^2\ (\text{both coordinates
continuous}).
\]
Then $\mathcal S=\{(0,0),(M,M)\}$, $\OPT=0$, $L_x=2$, $L_z=0$, and
$F_M(v)\ge\frac1{20}\dist(v,\mathcal S)^2$ on $X$, so $\kappa_S\le40$; one bag
$\{x,z\}$ covers both scopes ($p=2$). Let $\theta\in(0,1/4]$, $h\in(0,M]$ and
$c\in X$. Let $G_x$ be the graded grid (Definition~\ref{def:graded}) of
$[0,M]$ with center $c_x$, mesh $h$ and grading $\theta$, let $G_z\subseteq[0,M]$ be a finite set containing $0$
and $M$, and use the corrections $d_x=L_xw_x^2/8$ with $L_x\ge2$ and any
$d_z\ge0$. Then
\begin{equation}\label{eq:twocenters}
\beta=\min_{G_x\times G_z}Q\le-\max\Bigl\{\frac{\theta^2M^2}{379},\
\frac{h^2}{20}\Bigr\}.
\end{equation}
Consequently:
\begin{enumerate}[label=(\alph*)]
\item Every stage of TRIAL (Algorithm~\ref{alg:trial}), and hence every stage
of CT, satisfies~\eqref{eq:twocenters} with its grading $\theta$ and its mesh
$h=h_{xj}$ of coordinate $x$. Since $U\ge\OPT=0$, its certified gap $U-\beta$
is at least $\max\{\theta^2M^2/379,h^2/20\}$. The same holds for CT with a
common mesh (Lemma~\ref{lem:commonmesh}).
\item If such a stage has $\beta\ge-\varepsilon$, in particular if it
certifies a gap $U-\beta\le\varepsilon$, then its $x$-grid has more than
$M/(48\sqrt\varepsilon)$ nodes and its table for the bag $\{x,z\}$ has more
than $M/(24\sqrt\varepsilon)$ entries.
\item Any procedure that stops only after such a stage with $\beta>-1$, such
as EX (Algorithm~\ref{alg:ex}), builds a table with more than $M/24$ entries.
For integer $M$ the input length is $I=O(\log M)$, so this work is
$2^{\Omega(I)}$, although $\kappa_S\le40$ and $p=2$.
\end{enumerate}
\end{proposition}

\begin{proof}
\emph{The instance.} $F_M=x^2+z(M-2x)$. For $x<M/2$ the minimum over $z$ is
at $z=0$ with value $x^2$; for $x>M/2$ it is at $z=M$ with value $(M-x)^2$;
at $x=M/2$ the value is $M^2/4$. So $\mathcal S$ and $\OPT$ are as stated,
and the Hessian $\bigl(\begin{smallmatrix}2&-2\\-2&0\end{smallmatrix}\bigr)$
gives $L_x=2$, $L_z=0$. The map $(x,z)\mapsto(M-x,M-z)$ preserves $F_M$ (a
direct expansion) and swaps the minimizers, so for growth we may assume
$x\le M/2$. If $x\le M/4$, then $M-2x\ge M/2\ge z/2$ and
$F_M\ge x^2+z^2/2\ge\frac12\norm{(x,z)}^2$. If $M/4<x\le M/2$, then
$F_M\ge x^2\ge M^2/16$ while $\norm{(x,z)}^2\le5M^2/4$. In both cases
$F_M\ge\frac1{20}\dist((x,z),\mathcal S)^2$, and
$\kappa_S\le L/g_S=40$ with $g_S=1/20$.

\emph{The bound~\eqref{eq:twocenters}.} By symmetry assume $M-c_x\ge M/2$
(otherwise use the reflected argument with $z=0$ and the endpoint $0$). Let
$c_x=v_0<v_1<\dots<v_k=M$ be the $x$-nodes on the right of the center, so
$k\ge1$, and put $\delta_m=v_m-c_x$. All steps except the last are unclipped:
$\delta_{m+1}=(1+\theta)\delta_m+h$ for $m<k-1$, and
$\delta_k\le(1+\theta)\delta_{k-1}+h$. The node $M$ lies in $G_z$,
$F_M(v,M)=(M-v)^2$, every correction is nonnegative, and
$d_x(v)\ge w_x(v)^2/4$.

If $k=1$, the node $M$ has the adjacent interval $[c_x,M]$ of length at least
$M/2$, so $\beta\le Q(M,M)\le-M^2/16$, and
$M^2/16\ge\max\{\theta^2M^2/379,h^2/20\}$ since $\theta\le1/4$ and $h\le M$.

If $k\ge2$, put $\alpha=v_{k-1}-v_{k-2}$ and $\alpha'=M-v_{k-1}$. Then
$Q(M,M)\le-\alpha'^2/4$ and $Q(v_{k-1},M)\le\alpha'^2-\alpha^2/4$. If
$\alpha'^2\ge\alpha^2/5$, the first is at most $-\alpha^2/20$; otherwise the
second is below $-\alpha^2/20$; so $\beta\le-\alpha^2/20$. Now
$\alpha=h+\theta\delta_{k-2}\ge h$, and the recursion gives
$\delta_{k-2}\ge(\delta_k-(2+\theta)h)/(1+\theta)^2$. If $h\le M/16$, then
with $\delta_k\ge M/2$ and $\theta\le1/4$,
$\delta_{k-2}\ge(\frac M2-\frac{9M}{64})\cdot\frac{16}{25}=\frac{23M}{100}$, so
$\alpha\ge\frac{23}{100}\theta M$. If $h>M/16$, then
$\alpha\ge h>M/16\ge\theta M/4$. Hence $\alpha\ge\max\{h,\frac{23}{100}\theta M\}$
and $\beta\le-\max\{h^2,(0.23\,\theta M)^2\}/20$, which
implies~\eqref{eq:twocenters} because $0.23^2/20>1/379$.

\emph{Consequences.} (a) Since $L_z=0$, coordinate $z$ is not in $P$: TRIAL
uses $G_z=\{0,M\}$ and never filters $z$. Filtering with $U\ge\OPT$ keeps
both minimizers (Proposition~\ref{prop:filter}), so the $x$-interval of every
stage is $[0,M]$. At a stage $j\ge1$ the mesh is $h_{xj}=h_{x0}2^{-j}<M$ and
the center lies in $X$, so~\eqref{eq:twocenters} applies. At stage $0$ the
center is $\ell=(0,0)$ and $M\le h_{x0}<2M$ (Section~\ref{sec:growth}), so
$G_x=\{0,M\}$, $d_x(M)\ge M^2/4$ and $\beta\le Q(M,M)\le-M^2/4$; this is at
most $-\max\{\theta^2M^2/379,h_{x0}^2/20\}$ because $h_{x0}^2/20<M^2/5$.
With a common mesh, $h_j=M2^{-j}\le M$ at every stage.

(b) By~(a), $\beta\ge-\varepsilon$ gives $\theta M\le\sqrt{379\varepsilon}$
and $h\le\sqrt{20\varepsilon}$, so every step of the $x$-grid is at most
$h+\theta M\le(\sqrt{20}+\sqrt{379})\sqrt\varepsilon<24\sqrt\varepsilon$. The
side of $[0,M]$ from $c_x$ to the farther endpoint has length at least $M/2$,
so it takes at least $M/(48\sqrt\varepsilon)$ steps and contains more than
$M/(48\sqrt\varepsilon)$ nodes. Since $\abs{G_z}=2$, the table of the bag
$\{x,z\}$ has more than $M/(24\sqrt\varepsilon)$ entries.

(c) Apply (b) with $\varepsilon=1$. EX stops only when
$F(\tilde x)-\beta<1/(\Omega W)\le1$ for the bound $\beta$ of the successful
stage of CT, and $F(\tilde x)\ge\OPT=0$, so that stage has $\beta>-1$. The
data of $F_M$ are $1$, $-2$, $M$ and the endpoints $0$ and $M$, so
$I=O(\log M)$.
\end{proof}

Other methods solve this instance easily; it shows that the failure comes
from the single center, not from conditioning. Replacing hulls by unions of
retained cells, on uniform meshes, removes it, at the price of
$O(r\sqrt{n\kappa_S})$ instead of $O(\sqrt{\bar\kappa}\log(n+2))$ nodes per
coordinate.

\begin{algorithm}[UC: uniform cells]\label{alg:uc}
The input is an accuracy $\varepsilon>0$. For each coordinate keep a finite
set $\mathcal C_i$ of \emph{retained cells}, closed intervals of positive
length; initially $\mathcal C_i=\{[\ell_i,u_i]\}$. Let $J$ be the least
$j\ge0$ with $\frac12nLs^24^{-j}\le\varepsilon$, start with the incumbent
$\hat x=\ell$ and $U=F(\ell)$, and for $j=0,1,\dots,J$ with $h_j=s2^{-j}$:
\begin{enumerate}[label=(\roman*)]
\item \emph{Refine.} Let $\mathcal G_{ij}=\{\ell_i+kh_j:k\in\Z_{\ge0}\}$ for
$i\in I_C$ and $\mathcal G_{ij}=\{\ell_i+\lceil kh_j\rceil:k\in\Z_{\ge0}\}$ for
$i\in I_Z$. Split every retained cell $[a,a']$ at the points of
$\mathcal G_{ij}\cap(a,a')$; the resulting intervals are the \emph{stage
cells}, and $G_i$ is the set of their endpoints. Let $w_i(v)$ be the largest
effective width of a stage cell with endpoint $v$, and define $d_i$, $D$,
$Q$, $\beta_j$ and $m_i$ by~\eqref{eq:correction}--\eqref{eq:minmarginal} on
$G=\prod_iG_i$.
\item \emph{Solve.} Compute $\beta_j$, a minimizer $y^{(j)}$ of $Q$ over $G$
and all $m_i$ by Lemma~\ref{lem:dp}; if $F(y^{(j)})<U$, set $\hat x=y^{(j)}$
and $U=F(y^{(j)})$.
\item \emph{Filter.} Let $\mathcal C_i$ be the set of stage cells $[a,a']$
with $\min\{m_i(a),m_i(a')\}\le U$. Do not take hulls.
\end{enumerate}
Return $\hat x$, $U$ and $\beta_J$. The \emph{domain} of stage $j$ is
$X_j=X\cap\prod_i\bigcup\mathcal C_i$, with $\mathcal C_i$ as at the start of
stage $j$.
\end{algorithm}

Since $\mathcal G_{i,j-1}\subseteq\mathcal G_{ij}$ (for integers,
$\lceil kh_{j-1}\rceil=\lceil 2kh_j\rceil$), every stage cell is either a
segment between consecutive points of $\mathcal G_{ij}\cap[\ell_i,u_i]$ or
such a segment clipped at $u_i$. Hence a continuous stage cell has length at
most $h_j$; an integer stage cell has length at most $\lceil h_j\rceil$, and
if that length is at least two then $h_j>1$ and $\lceil h_j\rceil<2h_j$.
Consecutive points of $\mathcal G_{ij}$ are at least $h_j$ apart for
$i\in I_C$ and at least $\max\{1,\lfloor h_j\rfloor\}\ge h_j/2$ apart for
$i\in I_Z$. Each stage cell contains at most one point of
$\mathcal G_{i,j+1}$ in its interior.

\begin{lemma}[Uniform cells: bounds and safe filtering]\label{lem:cells}
Assume only Definition~\ref{def:curvature}. At every stage $j$ of UC,
$\beta_j\le\min_{X_j}F$, and $F(x)\ge\min\{m_i(a),m_i(a')\}$ for every
$x\in X_j$ and every stage cell $[a,a']$ of a coordinate $i$ with
$x_i\in[a,a']$. Every global minimizer lies in every $X_j$, and
$\beta_J\le\OPT\le U$.
\end{lemma}

\begin{proof}
The stage cells of coordinate $i$ form a finite family of intervals with
endpoints in $G_i$, every node is an endpoint of one of them, $w_i$ is
defined from this family, and $X_j$ is the intersection of $X$ with the
product of their unions. The first two claims are therefore
Proposition~\ref{prop:cellwise}(b) and~(c) in the form for families of
intervals (Section~\ref{sec:grids}). If step~(iii) of stage $j$ removes a
point $x\in X_j$, then for some $i$ every stage cell containing $x_i$ is
discarded, so $F(x)>U$ for the threshold $U$ used at stage $j$. As
$U\ge\OPT$, global minimizers are never removed; by induction every $X_j$
contains $\mathcal S$, and $\beta_J\le\min_{X_J}F=\OPT\le U$.
\end{proof}

\begin{theorem}[Uniform cells]\label{thm:cells}
Assume Definition~\ref{def:curvature}. For every $\varepsilon>0$,
UC$(\varepsilon)$ returns a feasible $\hat x$ with
$\beta_J\le\OPT\le F(\hat x)\le\beta_J+\varepsilon$, together with a
certificate (all stage tables and filtering decisions) whose validity uses
only Definition~\ref{def:curvature}. Suppose moreover that $\mathcal S$ is
finite, $r=\max_i\abs{\mathcal S_i}$, and $F$ has set growth with constant
$g_S$. Then every grid of every stage has at most
\begin{equation}\label{eq:cellcount}
K_S=12\,r\,\bigl(2\sqrt{n\kappa_S}+1\bigr)
\end{equation}
nodes per coordinate. For an explicit rational quadratic, UC$(2^{-q})$ uses
$O\bigl(p(\abs{\mathcal A}+n+N)K_S^p\bigr)\poly(I+q)$ bit operations, and
EX with UC in place of CT returns an exact minimizer in
$O\bigl(p(\abs{\mathcal A}+n+N)K_S^p\bigr)\poly(I+\log\kappa_S)$ bit
operations. None of $g_S$, $\mathcal S$ and $r$ is an input.
\end{theorem}

\begin{proof}
\emph{Gap.} Every stage cell of positive effective width has length at most
$h_j$ (continuous) or below $2h_j$ (integer), so
$D(y)\le\sum_iL_i(2h_j)^2/8\le\frac12nLh_j^2$ for every $y\in G$. Hence
$U-\beta_j\le F(y^{(j)})-\beta_j=D(y^{(j)})\le\frac12nLh_j^2$, which is at
most $\varepsilon$ at $j=J$. Validity is Lemma~\ref{lem:cells}.

\emph{Count.} Fix a stage $j$ and a cell $[a,a']$ retained by step~(iii).
One endpoint $v$ has $m_i(v)\le U$; let $z\in G$ attain $m_i(v)$, so $z_i=v$
and $Q(z)\le U$. By Lemma~\ref{lem:cells}, $\beta_j\le\OPT$, so
$U\le F(y^{(j)})=\beta_j+D(y^{(j)})\le\OPT+\frac12nLh_j^2$ and
\[
F(z)-\OPT=Q(z)+D(z)-\OPT\le U+\tfrac12nLh_j^2-\OPT\le nLh_j^2 .
\]
Set growth gives $\dist(z,\mathcal S)^2\le nLh_j^2/g_S\le n\kappa_Sh_j^2$, so
$\abs{v-\bar v}\le\sqrt{n\kappa_S}\,h_j$ for some $\bar v\in\mathcal S_i$.
The endpoints of stage cells are points of $\mathcal G_{ij}$ or $u_i$, and
consecutive points of $\mathcal G_{ij}$ are at least $h_j/2$ apart; so for
each $\bar v\in\mathcal S_i$ at most $4\sqrt{n\kappa_S}+2$ endpoints lie
within $\sqrt{n\kappa_S}\,h_j$ of $\bar v$. Each endpoint belongs to at most
two stage cells, so at most $2r(4\sqrt{n\kappa_S}+2)$ cells are retained.
Each of them contains at most one point of $\mathcal G_{i,j+1}$ in its
interior, so the grid of stage $j+1$ has at most
$3\cdot2r(4\sqrt{n\kappa_S}+2)=K_S$ nodes. Stage $0$ has two nodes per
coordinate.

\emph{Work and bit length.} By Lemma~\ref{lem:dp} each stage takes
$O(p(\abs{\mathcal A}+n+N)K_S^p)$ table operations, and
$J=O(\log(nLs^2/\varepsilon))=O(I+q)$. Continuous nodes have the form
$\ell_i+ks2^{-j}$ and integer nodes are integers, so all nodes have
denominators dividing $\Delta2^j$, with $\Delta$ as in
Section~\ref{sec:exact-data}, and bit length $O(I+j)$; factor values,
corrections and messages have bit length polynomial in $I+j$. For exact
output, Theorem~\ref{thm:transfer}(a) applies once $2^{-q}\le\varepsilon_S$,
and $\log_2(1/\varepsilon_S)\le\poly(I)+\log_2\kappa_S$ as in the proof of
Theorem~\ref{thm:exact} (if $L=0$, UC is exact at stage $0$). EX doubles $q$,
so it makes $O(\log(I+\log\kappa_S))$ calls, each with
$q\le\poly(I)+2\log_2\kappa_S$.
\end{proof}

\begin{remark}[What grading buys]\label{rem:grading}
With $r=1$ the minimizer is unique and $\kappa_S=\kappa$, and
Theorem~\ref{thm:cells} gives accuracy-independent grids with
$O(\sqrt{n\kappa})$ nodes per coordinate without gradings or caps. The graded
single-center grids of Section~\ref{sec:growth} need only
$O(\sqrt{\bar\kappa}\log(n+2))$ nodes per coordinate (Lemma~\ref{lem:states}),
which gives the $f(p,\bar\kappa)\poly(I+q)$ bound of
Theorem~\ref{thm:approx}; Proposition~\ref{prop:twocenters} shows that they
lose this bound when the minimizer is not unique. Whether exact output admits
an $f'(p,\kappa_S,r)\poly(I)$ bound when $\mathcal S$ is finite is open. When
$\mathcal S$ is a continuum, Proposition~\ref{lim:prop:setgrowth} shows that
runs of filtering with thresholds $U_j\ge\OPT$, in particular TRIAL and CT,
can need $1+1/(2\sqrt\varepsilon)$ nodes per coordinate to certify accuracy
$\varepsilon$; whether other certificates avoid this is open.
\end{remark}


\subsection{Coordinatewise concave quadratics: every optimizer at once}\label{sec:endpointset}

Write $E_i=\{\ell_i,u_i\}$ and $E(B)=\prod_{i\in B}E_i$. For $x\in\bar X$
and $v_i\in E_i$ put
\[
e_i(v_i;x_i)=\begin{cases}u_i-x_i,&v_i=\ell_i,\\ x_i-\ell_i,&v_i=u_i,\end{cases}
\qquad
\pi_t(v;x)=\prod_{i\in B_t}\frac{e_i(v_i;x_i)}{u_i-\ell_i}\quad(v\in E(B_t)).
\]
Thus $\pi_t(\cdot;x)$ is the law of $Y_{B_t}$ when every coordinate of $x$
is rounded independently to an endpoint with mean $x_i$. Root $T$ at $t_0$,
let $\mathrm{ch}(t)$ be the children of $t$, and write
$B_t^{\uparrow}=B_t\cap B_{\mathrm{pa}(t)}$, with $B_{t_0}^{\uparrow}=\emptyset$.
Let $q_t$ be the sum of the factors assigned to bag $t$.

\begin{lemma}[Endpoint Bellman residual identity]\label{lem:endpointid}
Assume $H_{ii}\le0$ for every $i$. Compute, in post-order and with all
arguments ranging over endpoint labels,
\begin{align*}
M_t(v_{B_t^\uparrow})&=\min_{v_{B_t\setminus B_t^\uparrow}}
\Bigl[q_t(v_{B_t})+\sum_{t'\in\mathrm{ch}(t)}M_{t'}(v_{B_{t'}^\uparrow})\Bigr],\\
r_t(v_{B_t})&=q_t(v_{B_t})+\sum_{t'\in\mathrm{ch}(t)}M_{t'}(v_{B_{t'}^\uparrow})
-M_t(v_{B_t^\uparrow})\ \ge0 .
\end{align*}
Then the scalar $M_{t_0}$ equals $\OPT$, it is attained at an endpoint
assignment, and for every $x\in\bar X$
\begin{equation}\label{eq:endpointid}
F(x)-\OPT=\sum_{t}\sum_{v\in E(B_t)}r_t(v)\,\pi_t(v;x)
+\frac12\sum_{i=1}^n(-H_{ii})(x_i-\ell_i)(u_i-x_i).
\end{equation}
All tables and residuals are computed in
$O\bigl(p\,2^p(N+|\mathcal A|)\bigr)$ table operations on rationals of
bit length polynomial in $I$.
\end{lemma}

\begin{proof}
Since $B_{t'}^{\uparrow}\subseteq B_t$ for a child $t'$, the bracket is a
function of $v_{B_t}$, and $r_t\ge0$ because $M_t$ is its minimum over the
labels of $B_t\setminus B_t^\uparrow$.

\emph{Telescoping.} Let $v\in E([n])$. Every nonroot table $M_t$ occurs in
$\sum_tr_t(v_{B_t})$ once with a plus sign, in the residual of its parent, and
once with a minus sign. Hence $\sum_tr_t(v_{B_t})=\sum_tq_t(v_{B_t})-M_{t_0}
=F(v)-M_{t_0}$.

\emph{Expectation.} Fix $x\in\bar X$ and let $Y$ be the independent endpoint
rounding with $\Pr(Y_i=v_i)=e_i(v_i;x_i)/(u_i-\ell_i)$. Then $\E Y_i=x_i$,
$\Var Y_i=(x_i-\ell_i)(u_i-x_i)$, $\E Y_iY_j=x_ix_j$ for $i\ne j$,
and $\E Y_i^2=x_i^2+\Var Y_i$. Therefore
$\E F(Y)=F(x)+\frac12\sum_iH_{ii}(x_i-\ell_i)(u_i-x_i)$, while
$\E r_t(Y_{B_t})=\sum_vr_t(v)\pi_t(v;x)$. Taking expectations in the
telescoping identity gives \eqref{eq:endpointid} with $M_{t_0}$ in place
of $\OPT$.

\emph{Value.} Every term on the right of \eqref{eq:endpointid} is
nonnegative on $\bar X$, so $F\ge M_{t_0}$ on $\bar X\supseteq X$. Trace back:
choose labels of $B_{t_0}$ attaining $M_{t_0}$; for each child $t'$ of a
processed bag $t$, choose the labels of $B_{t'}\setminus B_{t'}^\uparrow$
attaining $M_{t'}(v_{B_{t'}^\uparrow})$. By running intersection, a variable
belongs to $B_t\setminus B_t^\uparrow$ exactly for the bag $t$ nearest the
root among those containing it, and to $B_{t}^{\uparrow}$ for every other bag
containing it. Hence every variable receives exactly one label, and the
resulting endpoint assignment $v^*$ has all residuals zero. By telescoping
$F(v^*)=M_{t_0}$. Both endpoints of every $X_i$ are feasible, so $v^*\in X$
and $\OPT=M_{t_0}$.

\emph{Checking.} The proof uses only the nonnegativity of the residuals and
one endpoint assignment with zero residuals. A checker therefore needs the
tables $M_t$ and that assignment, not the minimization that produced them; any
tables with these two properties give \eqref{eq:endpointid}.

\emph{Cost.} A bag table has at most $2^p$ entries, and children are combined
by addition. By induction on the tree, every message entry is the sum of the
assigned factors over the subtree at one endpoint assignment. Its bit length
is therefore bounded by the bit length of a common denominator of all endpoint
monomial values plus $\log_2$ of the number of monomials times their maximum
magnitude, both polynomial in $I$.
\end{proof}

\begin{theorem}[All optimizers of a coordinatewise concave quadratic]\label{thm:endpointset}
Under the hypothesis of Lemma~\ref{lem:endpointid}, a point $x\in X$ is a
global optimizer if and only if
\begin{enumerate}
\item[(a)] $x_i\in\{\ell_i,u_i\}$ for every $i$ with $H_{ii}<0$, and
\item[(b)] $\prod_{i\in B_t}e_i(v_i;x_i)=0$ for every bag $t$ and every
$v\in E(B_t)$ with $r_t(v)>0$.
\end{enumerate}
These are at most $n+N2^p$ factored polynomial equations of degree at most
$p$, together with the original bounds and integrality. They are computed and
checked in $2^{O(p)}\poly(I)$ bit operations, without a growth assumption and
without enumerating interior integer labels or optimal components.
\end{theorem}

\begin{proof}
By \eqref{eq:endpointid}, $x$ is optimal exactly when every nonnegative term
vanishes. The diagonal terms vanish if and only if (a) holds, and
$\pi_t(v;x)=0$ exactly when one factor $e_i(v_i;x_i)$ is zero.
\end{proof}

Condition (b) is a support condition, not a condition on coordinate
projections. For $F=x+y-2xy=x(1-y)+y(1-x)$ on $[0,1]^2$, both optimal
corners $(0,0)$ and $(1,1)$ use every endpoint of every coordinate, yet the
two positive residuals give $x(1-y)=0$ and $y(1-x)=0$, which exclude every
other point. A sum of $m$ disjoint copies has $2^m$ isolated optimizers and
a description with $2m$ equations.

The equations have a finite combinatorial form. For $x\in\bar X$ define its
\emph{face pattern} $\chi(x)\in\{\mathsf L,\mathsf U,\ast\}^n$ by
$\chi_i(x)=\mathsf L$ if $x_i=\ell_i$, $\chi_i(x)=\mathsf U$ if $x_i=u_i$, and
$\chi_i(x)=\ast$ otherwise. For a pattern $\chi$ on $B\subseteq[n]$ let
$\operatorname{cube}(\chi)\subseteq E(B)$ consist of the $v$ with $v_i=\ell_i$
where $\chi_i=\mathsf L$ and $v_i=u_i$ where $\chi_i=\mathsf U$. For a
pattern $\chi$ on $[n]$, the closed face
$\operatorname{face}(\chi)=\{x\in\bar X:x_i=\ell_i\text{ if }\chi_i=\mathsf L,\
x_i=u_i\text{ if }\chi_i=\mathsf U\}$ of $\bar X$ has the vertex set
$\operatorname{cube}(\chi)$. Then $\pi_t(v;x)>0$ if and only if
$v\in\operatorname{cube}(\chi(x)_{B_t})$. Let $\Sigma_i=\{\mathsf L,\mathsf U\}$,
with $\ast$ added when $H_{ii}=0$ and $X_i$ has a point strictly between
$\ell_i$ and $u_i$. Call $\chi\in\prod_i\Sigma_i$ \emph{admissible} if, for
every bag, $\operatorname{cube}(\chi_{B_t})\subseteq Z_t:=\{v\in E(B_t):r_t(v)=0\}$.

\begin{corollary}[Optimal faces form a tree-structured constraint problem]\label{cor:facecsp}
Under the hypothesis of Lemma~\ref{lem:endpointid}:
\begin{enumerate}
\item[(i)] $x\in X$ is optimal if and only if $\chi(x)\in\prod_i\Sigma_i$
and $\chi(x)$ is admissible.
\item[(ii)] Replacing any $\ast$ of an admissible pattern by $\mathsf L$ or
$\mathsf U$ keeps it admissible, and
$\mathcal S=\bigcup_{\chi\text{ admissible}}\bigl(X\cap\operatorname{face}(\chi)\bigr)$.
Distinct admissible patterns give disjoint nonempty sets
$\{x\in X:\chi(x)=\chi\}$.
\item[(iii)] Admissible patterns are the solutions of a constraint problem on
the same tree decomposition, with at most three labels per variable and bag
relations
$\mathrm{Adm}_t=\{\chi\in\prod_{i\in B_t}\Sigma_i:\operatorname{cube}(\chi)
\subseteq Z_t\}$. Consequently, in $2^{O(p)}(N+n)\poly(I)$ bit operations
one can decide whether the optimizer is unique, compute the dimension of
$\mathcal S$, count the admissible patterns, count $|\mathcal S|$ when
$\mathcal S$ is finite, compute $\dist(y,\mathcal S)^2$ and a nearest
optimizer for a rational $y$, and minimize a linear function over
$\mathcal S$.
\end{enumerate}
\end{corollary}

\begin{proof}
(i) Condition (a) of Theorem~\ref{thm:endpointset} says $\chi_i(x)\ne\ast$
when $H_{ii}<0$; and $\chi_i(x)=\ast$ is possible only if $X_i$ has a point
strictly between its endpoints. Condition (b) says that every $v$ with
$\pi_t(v;x)>0$ has $r_t(v)=0$, that is,
$\operatorname{cube}(\chi(x)_{B_t})\subseteq Z_t$.

(ii) Replacing $\ast$ by an endpoint shrinks every cube. If $\chi$ is
admissible and $x\in X\cap\operatorname{face}(\chi)$, then $\chi(x)$ is
obtained from $\chi$ by replacing some $\ast$ entries by endpoints, so $x$ is
optimal by (i). Conversely $x\in\mathcal S$ lies in
$\operatorname{face}(\chi(x))$. A pattern $\chi$ determines the set
$\{x\in X:\chi(x)=\chi\}$, which is nonempty by the definition of $\Sigma_i$.

(iii) Each $\mathrm{Adm}_t$ is computed by testing at most $3^p$ patterns
against at most $2^p$ residual entries. A constraint problem whose relations
live on the bags of a tree decomposition is solved by nonserial dynamic
programming over $(T,(B_t))$ in the min-sum, max-sum and sum-product
semirings, with $O(p\,3^p)$ operations per bag \cite{Dechter1999}; each unary
cost is charged once, at the bag nearest the root that contains its variable.
The optimizer is unique if and only if exactly one pattern is admissible and
it has no $\ast$, by the disjointness in (ii). The dimension of $\mathcal S$
is the maximum number of $\ast$ entries on continuous coordinates. When
$\mathcal S$ is finite,
$|\mathcal S|=\sum_\chi\prod_{i:\chi_i=\ast}|X_i\cap(\ell_i,u_i)|$. By (ii),
$\dist(y,\mathcal S)^2$ is the minimum over admissible $\chi$ of the sum of
the unary costs $(y_i-\ell_i)^2$, $(y_i-u_i)^2$ and $\dist(y_i,X_i)^2$ for
$\chi_i=\mathsf L$, $\mathsf U$ and $\ast$. A linear objective attains its
minimum over each $X\cap\operatorname{face}(\chi)$ at an endpoint pattern, so
it is minimized by the same recursion over $\{\mathsf L,\mathsf U\}$ labels.
All numbers are rationals of polynomial bit length.
\end{proof}

\begin{remark}[Two-label integer coordinates and multilinear objectives]\label{rem:endpointext}
If $X_i=\{\ell_i,\ell_i+1\}$, replacing $\frac12H_{ii}x_i^2$ by
$\frac12H_{ii}\bigl((2\ell_i+1)x_i-\ell_i(\ell_i+1)\bigr)$ does not change $F$
on $X$. Hence the sign hypothesis is needed only for coordinates with at least
three feasible values; all binary quadratic programs are covered. The proofs
also apply verbatim to $F=\sum_af_a+\sum_i\psi_i(x_i)$ in which each $f_a$ is
affine in each of its variables separately and each $\psi_i$ is concave on
$[\ell_i,u_i]$. Then $\E f_a(Y)=f_a(x)$, and the diagonal term of
\eqref{eq:endpointid} becomes $\psi_i(x_i)$ minus its chord at $x_i$. That
difference vanishes at an interior point exactly when $\psi_i$ is affine on
$[\ell_i,u_i]$, which replaces the condition $H_{ii}<0$ in (a).
\end{remark}

Endpoint optimality for nonpositive diagonals is classical. For multilinear
functions on a box, Rosenberg characterized the optimal set as the union of
the faces all of whose vertices are optimal \cite{Rosenberg1972}. For a
quadratic on a continuous box with $H_{ii}=0$ for all $i$, this is
Corollary~\ref{cor:facecsp}(ii), because by (i) a pattern is admissible
exactly when all vertices of its face are optimal. Del Pia and
Khajavirad use endpoint optimality to separate binary-valued from continuous
variables \cite{DelPiaKhajavirad2026}. Min-sum messages are standard
\cite{Dechter1999}, and nonnegative reparameterizations whose zero entries
describe all optimal labelings of a discrete problem are known from max-sum
inference \cite{WainwrightJaakkolaWillsky2005,Werner2007}.
Theorem~\ref{thm:endpointset} combines these facts and extends them to
strictly concave coordinates and mixed boxes. The new part is the factored
form on the decomposition: Corollary~\ref{cor:facecsp} turns the optimal set
into a tree-structured pattern problem, which gives uniqueness, dimension,
counts, distances and linear optimization over $\mathcal S$ with $2^{O(p)}$
operations per bag.

\subsection{The diagonal Lagrangian certificate and unknown growth}\label{sec:diagcert}

In this subsection $X=\bar X$ is a continuous box and $F$ is a rational
quadratic. A point $\bar x\in X$ is a \emph{box KKT point} if
$\partial_iF(\bar x)\ge0$ when $\bar x_i=\ell_i$, $\partial_iF(\bar x)\le0$
when $\bar x_i=u_i$, and $\partial_iF(\bar x)=0$ otherwise; every minimizer
is one. For such a point put
\[
\lambda_i(\bar x)=\frac{2\,|\partial_iF(\bar x)|}{u_i-\ell_i},\qquad
\Lambda(\bar x)=\diag(\lambda_1(\bar x),\ldots,\lambda_n(\bar x)).
\]
For $\lambda\in\R^n_{\ge0}$ define the Lagrangian dual function
\[
\Psi(\lambda)=\inf_{x\in\R^n}\Bigl[F(x)-\frac12\sum_i\lambda_i(x_i-\ell_i)(u_i-x_i)\Bigr].
\]

\begin{lemma}[Diagonal certificate, full optimal set and invariance]\label{lem:diagcert}
Let $X$ be a continuous box and $F$ a quadratic.
\begin{enumerate}
\item[(i)] For every box KKT point $\bar x$ and every $x\in\R^n$,
\begin{equation}\label{eq:diagid}
\begin{split}
F(x)-F(\bar x)={}&\frac12(x-\bar x)^{\T}\bigl(H+\Lambda(\bar x)\bigr)(x-\bar x)\\
&+\frac12\sum_i\lambda_i(\bar x)(x_i-\ell_i)(u_i-x_i).
\end{split}
\end{equation}
\item[(ii)] If $\bar x$ is a box KKT point and $H+\Lambda(\bar x)\succeq0$,
then $\bar x\in\mathcal S$ and
\begin{equation}\label{eq:diagset}
\mathcal S=\bigl\{x\in X:\ (H+\Lambda(\bar x))(x-\bar x)=0,\ \
\lambda_i(\bar x)(x_i-\ell_i)(u_i-x_i)=0\ \text{for all }i\bigr\}.
\end{equation}
\item[(iii)] $\Psi(\lambda)\le\OPT$ for every $\lambda\ge0$, and the
following are equivalent: (a) $H+\Lambda(\bar x)\succeq0$ for some
$\bar x\in\mathcal S$; (b) $H+\Lambda(\bar x)\succeq0$ for every
$\bar x\in\mathcal S$; (c) $\Psi(\lambda)=\OPT$ for some $\lambda\ge0$. When
they hold, the vector $\lambda(\bar x)$ is the same for every
$\bar x\in\mathcal S$, and it is the unique maximizer of $\Psi$.
\end{enumerate}
The proof is in Appendix~\ref{app:proximal}.
\end{lemma}

\begin{remark}[The diagonal certificate class]\label{rem:shor}
Call the instances satisfying (a)--(c) the \emph{diagonal certificate class}
$\mathfrak D$. Convex box QPs belong to it, since $\Lambda(\bar x)\succeq0$.
It also contains indefinite instances with tilted and disconnected optimal
sets. For example $F=(x-y+t/2)^2+\frac18t(1-t)$ on $[0,1]^3$ has diagonal
$(2,2,\frac14)$, the two optimal segments $\{(v,v,0):0\le v\le1\}$ and
$\{(v,v+\frac12,1):0\le v\le\frac12\}$, the multiplier
$\Lambda=\diag(0,0,\frac14)$ at every optimizer, and
$H+\Lambda=2(1,-1,\tfrac12)^{\T}(1,-1,\tfrac12)\succeq0$. The function
$\Psi$ is the Lagrangian dual for the constraints
$(x_i-\ell_i)(u_i-x_i)\ge0$, so condition (c) says that this dual,
equivalently the basic Shor semidefinite relaxation, is exact. Sufficient
global optimality conditions of this box-multiplier form are classical
\cite{JeyakumarRubinovWu2006,LiWuQuan2015}, and stronger relaxations with
complete exactness characterizations are studied in \cite{QiuYildirim2024}.
Lemma~\ref{lem:diagcert} is therefore not a new certificate. What we use is
its invariance: whichever optimizer a search finds, the same rational test
accepts it.
\end{remark}

Lemma~\ref{lem:diagcert} checks a candidate but does not produce one.
Candidates come from proximal stages: graded grids centered at the previous
output, with a proximal term instead of filtering, and a guess $\hat\kappa$
in place of the unknown $\kappa_S$.

\begin{algorithm}[PROX: proximal stages with guess $\hat\kappa$]\label{alg:prox}
The input is a rational $\hat\kappa\ge1$. Let $\mu\ge2$ be the least integer
with $4^{-\mu}\hat\kappa\le1/8$, and put
\[
\theta=2^{-\mu},\qquad\eta=\frac{L\theta^2}4,\qquad
\varrho=2\bigl\lceil\sqrt{\hat\kappa n}\,\bigr\rceil .
\]
A \emph{proximal stage} with center $c\in X$ and mesh $h>0$ forms the box
$X'=X\cap\prod_i[c_i-\varrho h,c_i+\varrho h]$, takes for $G_i$ the graded
grid of $X'_i$ with center $c_i$, mesh $h$ and grading $\theta$, computes the
corrections $d_i$ and $D$ of~\eqref{eq:correction} on $G=\prod_iG_i$, and
returns a minimizer $y^+$ of
\[
Q_\eta(y)=F(y)-D(y)+\eta\norm{y-c}^2\qquad(y\in G),
\]
computed by Lemma~\ref{lem:dp}; the correction and proximal terms are unary,
so the decomposition is unchanged. PROX$(\hat\kappa)$ runs proximal stages
$j=0,1,\dots$ with meshes $h_j=s2^{-j}$ and centers $c^{(0)}=\ell$ and
$c^{(j+1)}=y^{(j)}$, where $y^{(j)}$ is the output of stage $j$. There is no
filtering: every box is formed from $X$.
\end{algorithm}

\begin{lemma}[Proximal stages]\label{lem:proximal}
Let $K_\theta=8\theta^{-1}\lceil\log_2(n+2)\rceil$. In PROX$(\hat\kappa)$,
$\theta\le1/4$, $8\hat\kappa\theta^2\le1$ and $\theta^{-1}\le6\sqrt{\hat\kappa}$,
and:
\begin{enumerate}[label=(\roman*)]
\item every grid has at most $K_\theta$ nodes per coordinate;
\item if $F$ has set growth with a constant $g_S$ such that
$\kappa_S\le\hat\kappa$, then a proximal stage whose center satisfies
$\dist(c,\mathcal S)^2\le4\hat\kappa nh^2$ returns $y^+$ with
$F(y^+)-\OPT\le\frac{99}{256}Lnh^2$ and
$\dist(y^+,\mathcal S)^2\le\frac{99}{256}\hat\kappa nh^2$;
\item under the hypothesis of (ii) on $g_S$,
$\dist(y^{(j)},\mathcal S)^2\le\frac{99}{256}\hat\kappa nh_j^2$ for every
$j\ge0$;
\item stage $j$ costs $O\bigl(p(\abs{\mathcal A}+n+N)K_\theta^p\bigr)$ table
operations on rationals of bit length $O(I+j+\mu K_\theta)$.
\end{enumerate}
The proof is in Appendix~\ref{app:proximal}.
\end{lemma}

\begin{algorithm}[DISC: discovery under unknown growth]\label{alg:disc}
For $\hat\kappa=1,2,4,\dots$: let $J_{\hat\kappa}$ be the least $j\ge0$ with
$4^j\tau^2\ge4\hat\kappa ns^2$ ($\tau$ as in REC, Algorithm~\ref{alg:rec});
run PROX$(\hat\kappa)$ for stages $0,\dots,J_{\hat\kappa}$ and let
$y=y^{(J_{\hat\kappa})}$. Run REC$(y)$; if it does not fail, let $\bar x$ be
a vertex of the polytope of solutions of its linear system, as returned by
exact linear programming. If $\bar x$ is a box KKT point and
$H+\Lambda(\bar x)\succeq0$, stop and output $\bar x$, $\OPT=F(\bar x)$,
$\Lambda(\bar x)$ and the description~\eqref{eq:diagset} of $\mathcal S$;
otherwise continue with the next guess.
\end{algorithm}

\begin{theorem}[Discovery of a certified optimal set under unknown growth]\label{thm:diagdiscovery}
Let $X$ be a continuous box and $F$ a rational quadratic.
\begin{enumerate}[label=(\alph*)]
\item Every output of DISC is correct, with no assumption on the instance.
If $F\notin\mathfrak D$, DISC does not stop.
\item If $F\in\mathfrak D$, then $F$ has set growth with some constant
$g_S>0$, and for every such $g_S$, DISC stops no later than at the first
guess $\hat\kappa\ge\kappa_S$, so with $\hat\kappa\le2\kappa_S$.
\item Under the hypotheses of (b), DISC uses at most $f'(p,\kappa_S)\poly(I)$
bit operations, where $f'$ is a computable function that is nondecreasing
and polynomial in $\kappa_S$ for fixed $p$, and the polynomial in $I$ has an
absolute exponent. The output has bit length $\poly(I)$.
\end{enumerate}
Neither $\mathcal S$, nor $g_S$, nor a bound on $\kappa_S$ is supplied.
\end{theorem}

The proof is in Appendix~\ref{app:proximal}. An implementation needs a work
limit, because DISC does not stop outside $\mathfrak D$. An unsuccessful
guess certifies nothing about membership: with a guess $\hat\kappa<\kappa_S$,
PROX can return the same nonoptimal point at every stage within the budget
(Remark~\ref{rem:falseguess}). For fixed $p$ the bound is polynomial in
$\kappa_S$ and $I$: by Lemma~\ref{lem:proximal}, for every guess
$\hat\kappa\le2\kappa_S$ a bag table has at most
$K_\theta^p\le2(48p\sqrt{\hat\kappa})^p(n+2)\le2(68p\sqrt{\kappa_S})^p(n+2)$
entries, using~\eqref{eq:logabsorb} with $n$ in place of $n_P$, and the bit
lengths are polynomial in $I$ and $\sqrt{\kappa_S}\log(2\kappa_S)$.

\begin{proposition}[Discovery within $\mathfrak D$ is hard without the growth parameter]\label{prop:sshard}
For integers $a_0,a_1,\ldots,a_m$ with $|a_0|\le A=\sum_{i=1}^m|a_i|$, let
$\xi_0=0$ and
\[
F(x,\xi)=\sum_{i=1}^m(\xi_i-\xi_{i-1}-a_ix_i)^2+(\xi_m-a_0)^2
+\sum_{i=1}^mx_i(1-x_i)
\]
on $x\in[0,1]^m$, $\xi\in[-A,A]^m$. The bags $\{\xi_1,x_1\}$ and
$\{\xi_{i-1},\xi_i,x_i\}$ $(2\le i\le m)$ form a path decomposition with bag
size at most $3$. If some $\hat x\in\{0,1\}^m$ has $\sum_{i=1}^ma_i\hat x_i=a_0$, then
$F\in\mathfrak D$, $\OPT=0$, and $\mathcal S$ consists exactly of the points
$(\hat x,\hat\xi)$ with $\hat x$ such a solution and $\hat\xi$ its partial
sums. Otherwise $\OPT>0$. Consequently, unless $\mathrm P=\mathrm{NP}$:
\begin{enumerate}[label=(\roman*)]
\item no polynomial-time algorithm returns an exact optimizer for every
instance of $\mathfrak D$ with bag size three, even though $\OPT=0$ is known
in advance for these instances;
\item there is no polynomial $\pi$ such that every instance of this form in
$\mathfrak D$ has a set-growth constant with $\kappa_S\le\pi(I)$.
\end{enumerate}
\end{proposition}

The objective is the Subset Sum reduction of
\cite[Remark~2]{DelPiaKhajavirad2026} with unscaled partial sums; what the
proposition adds is that its feasible instances belong to $\mathfrak D$. The
proof is in Appendix~\ref{app:proximal}.

The optimal value of an instance of $\mathfrak D$ equals $\max\Psi$, the
value of a convex relaxation. Proposition~\ref{prop:sshard} shows that the
difficulty is in locating an optimizer, and that a parameter such as
$\kappa_S$ is needed for an efficient exact search.
Corollary~\ref{cor:facecsp} has no counterpart here. With $a_0=0$ and
integers $a_1,\ldots,a_m$ of both signs, the origin is an optimizer of an
instance of $\mathfrak D$, and it is the unique optimizer exactly when no
nonempty subset of $a_1,\ldots,a_m$ sums to zero, which is an NP-complete
question. So even with one optimizer and the full
description~\eqref{eq:diagset} in hand, deciding uniqueness is hard.
